\documentclass[11pt]{article}

\usepackage[margin=0.6in]{geometry}
\usepackage[T1]{fontenc}
\usepackage[utf8]{inputenc}
\usepackage{lmodern}
\usepackage{array}
\usepackage{booktabs}
\usepackage{longtable}
\usepackage{hyperref}
\usepackage{parskip}
\usepackage{ragged2e}
\usepackage{setspace}
\usepackage{enumitem}
\usepackage{needspace}
\usepackage[table]{xcolor}
\usepackage[natbibapa]{apacite}

\definecolor{PrimaryRed}{HTML}{8B1E1E}
\definecolor{SoftBlue}{HTML}{DDEEFF}
\definecolor{SoftGray}{HTML}{F4F4F4}

\hypersetup{colorlinks=true,linkcolor=black,urlcolor=black,citecolor=black}
\setlength{\parindent}{0pt}
\setstretch{1.08}
\setlist[itemize]{leftmargin=0.23in,itemsep=0.02in,topsep=0.04in}
\setlist[enumerate]{leftmargin=0.27in,itemsep=0.02in,topsep=0.04in}
\bibliographystyle{apacite}

\newcommand{\reviewsection}[1]{%
  \Needspace{4\baselineskip}
  \vspace{0.07in}\noindent{\normalsize\bfseries\color{PrimaryRed}#1}
  \par\vspace{0.02in}
}

\newcommand{\ptrtableheader}{%
  \rowcolor{PrimaryRed}
  \color{white}\bfseries Checklist Area &
  \color{white}\bfseries Completed Response \\
}

\begin{document}

\hypersetup{pageanchor=false}
\begin{titlepage}
\newgeometry{margin=0.7in}
\thispagestyle{empty}

\begin{center}
{\large University of Rochester\\}
{\normalsize Warner School of Education\\[0.15in]}

{\color{PrimaryRed}\rule{\textwidth}{1.4pt}}\\[0.12in]
{\Huge\bfseries From Compliance\\[0.07in]
to Communication}\\[0.15in]
{\huge Prevent--Teach--Reinforce Planning\\[0.07in]
Through Access, Agency, and\\[0.07in]
Richer Evidence}\\[0.18in]
{\Large EDE 448 --- Module 2 Assignment}\\[0.12in]

\vfill

{\color{PrimaryRed}\rule{0.78\textwidth}{0.7pt}}\\[0.22in]

\colorbox{SoftBlue}{%
\parbox{0.88\textwidth}{%
\centering\small\vspace{0.02in}
\textbf{Behavior as Evidence, Support as Access}\\[0.05in]
A case-based completion of the PTR Functional Behavioral Assessment Checklist using a non-identifying 2025--2026 teaching-placement event and a dignity-centered plan for future classroom practice.
}%
}
\end{center}

\vfill

\begin{center}
\renewcommand{\arraystretch}{1.18}
\setlength{\tabcolsep}{4pt}
\begin{tabular}{>{\bfseries\RaggedLeft}p{1.65in} p{5.0in}}
Student & Piter Zacari Garcia Bautista \\
Assignment & Prevent Teach Reinforce Planning \\
Course & EDE 448 --- Communication and Positive Behavioral Supports for Students with Autism and Other Complex Needs \\
Selected Resource & Appendix 4.1, PTR Functional Behavioral Assessment Checklist \\
Classroom Context & Pine Brook elementary technology/IGNITE placement, 2025--2026 \\
Interpretive Lenses & Puzzle Plan and EQUITAS \\
Due Date & July 31, 2026, 11:59 PM EDT \\
\end{tabular}
\end{center}

\vfill

\begin{center}
\begin{minipage}{0.92\textwidth}
\centering\itshape
Positive support should leave the learner with more communication, access, and agency than before the intervention.\\
{\color{PrimaryRed}\rule{4in}{0.8pt}}
\end{minipage}
\end{center}

\vfill

\begin{center}
\begin{minipage}{0.95\textwidth}
\small
This paper uses Prevent--Teach--Reinforce to examine a transition in which a learner protested the removal of a familiar laptop and then participated when the laptop and hands-on materials were allowed to coexist. The analysis separates observation from inference, treats the learner's response to support as evidence, and identifies what additional student, family, team, and repeated-observation data would be needed before calling the hypothesis a completed functional assessment.
\end{minipage}
\end{center}

\restoregeometry
\end{titlepage}
\hypersetup{pageanchor=true}
\pagenumbering{arabic}

\reviewsection{Why I Selected the PTR Functional Behavioral Assessment Checklist}

For this assignment, I selected Appendix 4.1, the PTR Functional Behavioral Assessment Checklist, from the second edition of \textit{Prevent--Teach--Reinforce} \citep{dunlap2019ptr}. The checklist supports the purpose of this assignment because it requires educators to examine the conditions surrounding behavior rather than beginning with a moral label or an adult-selected consequence. It organizes inquiry through three connected components: what can be changed to prevent unnecessary difficulty, what functionally useful skill or communication route should be taught, and what response will make that route effective for the learner.

This structure fits my understanding of behavior as communication, but it also makes that principle more accountable. Saying that behavior communicates something is not enough. I must define what I observed, distinguish observation from inference, examine the task and environment, ask whose knowledge is missing, test a tentative hypothesis, and monitor whether the support increases access and quality of life.

\reviewsection{Selected Classroom Context and Evidence Boundary}

I completed the checklist companion using one documented, non-identifying event from my 2025--2026 Pine Brook elementary technology/IGNITE placement. During a transition from a familiar laptop-based activity to building with blocks, a learner raised their voice, protested the laptop's removal, and continued holding or using the laptop when an adult attempted to move it away. I asked that the learner be allowed to keep the laptop while I invited them into the blocks activity. The learner began building and remained connected.

The event matters because the learner's participation changed after an adult changed the access condition. It does not establish a diagnosis, a fixed behavioral function, a time-of-day pattern, or the learner's internal state. One incident supports a useful functional hypothesis, not a completed functional behavioral assessment. I therefore marked unsupported checklist areas as unknown instead of inventing certainty.

\reviewsection{Observable Definition and Tentative Functional Hypothesis}

\textbf{Observable behavior.} During an attempted transition from laptop use to blocks, the learner raised their voice, verbally or physically protested removal of the laptop, and retained the device. This description avoids labels such as \textit{defiant}, \textit{noncompliant}, or \textit{disruptive}. It names what was observable and the condition immediately surrounding it.

\textbf{Tentative hypothesis.} When an adult abruptly attempted to remove a familiar laptop and begin a less familiar hands-on task, the learner protested and retained the device. The behavior may have delayed the abrupt transition and preserved access to a tool supporting predictability, regulation, confidence, communication, preferred activity, or several of these functions. When the laptop and blocks were allowed to coexist, the learner joined the activity.

\reviewsection{Completed PTR FBA Checklist Companion: Prevent}

\small
\setlength{\tabcolsep}{4pt}
\renewcommand{\arraystretch}{1.16}
\begin{longtable}{>{\RaggedRight\arraybackslash}p{1.62in} >{\RaggedRight\arraybackslash}p{5.42in}}
\ptrtableheader
\endfirsthead
\ptrtableheader
\endhead
Times most or least likely &
Unknown from the available record. One event cannot establish a time-of-day pattern. \\
\rowcolor{SoftGray}
Activity most associated &
Transition from familiar laptop use to a hands-on blocks activity while the device was being removed. \\
Activity least associated &
Unknown. Participation increased when laptop access and blocks were allowed to coexist, but this response to support does not establish a stable low-risk routine. \\
\rowcolor{SoftGray}
People associated with occurrence or nonoccurrence &
No reliable person pattern is established. The record shows different adult responses, not a conclusion about the learner's relationship with either adult. \\
Immediate circumstances &
Abrupt transition, attempted removal of a familiar item, and uncertainty about entering a different activity. \\
\rowcolor{SoftGray}
Academic skill match &
There is no evidence that the learner lacked the skill to build with blocks. The learner participated after the access condition changed. \\
Physical environment &
No specific noise, crowding, lighting, temperature, or sensory condition was documented. These factors remain unknown. \\
\rowcolor{SoftGray}
Events outside school &
Unknown. No family, health, sleep, medication, food, transportation, or prior-event information should be inferred. \\
Additional context &
The laptop was a familiar and apparently meaningful tool. Allowing it to remain removed an unnecessary either-or condition and supported entry into the shared activity. \\
\end{longtable}
\normalsize

\reviewsection{Completed PTR FBA Checklist Companion: Teach}

\small
\begin{longtable}{>{\RaggedRight\arraybackslash}p{1.62in} >{\RaggedRight\arraybackslash}p{5.42in}}
\ptrtableheader
\endfirsthead
\ptrtableheader
\endhead
Gain peer attention &
No evidence from the documented event. \\
\rowcolor{SoftGray}
Gain adult attention &
Adult attention followed the protest, but sequence alone does not establish attention as the behavioral function. \\
Obtain an item or activity &
Plausible. The learner acted to maintain laptop access. ``Preferred item'' may be incomplete because the laptop could also have supported regulation, predictability, competence, continuity, or communication. \\
\rowcolor{SoftGray}
Avoid or delay a transition &
Plausible. The event occurred during an attempted transition away from the laptop. \\
Avoid or delay the task &
Possible but not established. Participation in the blocks task increased when the laptop could remain, which argues against assuming rejection of the task itself. \\
\rowcolor{SoftGray}
Useful skills to teach or strengthen &
Communicate ``not yet,'' ``more time,'' ``keep this with me,'' ``show me the first step,'' ``help,'' or ``break''; choose a transition route; begin with a familiar support present; and re-enter through more than one participation mode. \\
\end{longtable}
\normalsize

\reviewsection{Completed PTR FBA Checklist Companion: Reinforce}

\small
\begin{longtable}{>{\RaggedRight\arraybackslash}p{1.62in} >{\RaggedRight\arraybackslash}p{5.42in}}
\ptrtableheader
\endfirsthead
\ptrtableheader
\endhead
What followed the behavior &
The forced removal stopped. The learner kept the laptop, received an invitation to build, and retained access to both materials. \\
\rowcolor{SoftGray}
What happened next &
The learner began building with blocks and continued participating. \\
Praise preference &
Unknown. Public praise, adult attention, or a specific reward should not be assumed to be reinforcing. \\
\rowcolor{SoftGray}
Acknowledgment pattern &
Not established by the single event. Repeated observation would be needed. \\
Correction pattern &
One attempted removal occurred before the support changed. A broader pattern is unknown. \\
\rowcolor{SoftGray}
Demonstrated interest or support &
Laptop access was meaningful in this event. It should be treated as a contextual support, not a permanent or universal reinforcer. \\
Natural reinforcement &
The learner retained continuity and meaningful choice while gaining access to the hands-on activity and shared participation. \\
\end{longtable}
\normalsize

\reviewsection{Prevent: Change the Conditions Before Escalation}

The prevent component helped me locate the concern in a specific transition instead of locating it entirely inside the student. The immediate conditions included an expected change from a familiar device to a different material, attempted removal of the familiar item, and a less familiar entry point. Prevention could include:

\begin{itemize}
  \item previewing the transition with a visible first/next sequence;
  \item showing and modeling the hands-on task before requesting a switch;
  \item offering a short overlap period in which both tools remain available;
  \item providing choices such as laptop beside blocks, observe-first, partner start, or a learner-selected timer;
  \item giving one direction at a time and making the entry step visible; and
  \item avoiding removal of a familiar or regulating tool before understanding its role.
\end{itemize}

These supports are not rewards for protest. They make the transition understandable before the learner must intensify communication. My EDU486 camp work reinforces this point. Youth asked for short explanations, visible sequences, and opportunities to begin quickly. Our field reflection also showed that rushed closure, inconsistently visible response options, and a timed count could introduce barriers. Prevention is not only an individualized accommodation. It is also the quality of the instructional design.

\reviewsection{Teach: Build a Communication Route That Works}

The teach component asks what the learner could do that would meet the same need more safely, clearly, or efficiently. Useful messages in this case include ``not yet,'' ``more time,'' ``keep this with me,'' ``show me,'' ``help,'' and ``break.'' Those messages should be available through speech, text, visuals, gesture, AAC, pointing, or partner-supported communication.

\citet{jorgensen2018competence} argues that presuming competence is necessary but insufficient; educators must construct communicative competence by making communication available across the day. If a student lacks a reliable way to negotiate a transition, punishing the current signal will not create a better one. The team must teach and honor a route that works during actual classroom pressure. Adults therefore need instruction too: they must learn to recognize protest as information, make the communication option available, and respond consistently whenever safety permits.

\reviewsection{Reinforce: Strengthen Agency Rather Than Compliance}

The reinforce component helped me distinguish reinforcement from bribery or adult-selected praise. In the documented event, the natural result of the support was continued access to a familiar tool and successful entry into a shared activity. Future reinforcement should make the replacement communication effective. If the learner uses ``more time'' or ``keep this with me,'' the adult should respond consistently when safely possible. The plan can also recognize joining, re-entering, trying another material, communicating a boundary, or contributing through an alternate route.

I would not define success as becoming quiet, surrendering the laptop, making eye contact, or transitioning only on the adult's preferred schedule. First-hand accounts collected by \citet{gardner2017firsthand} warn that behavioral intervention can produce masking, shame, trauma, and weakened boundaries when appearing typical becomes the real goal. Reinforcement should increase communication, self-advocacy, participation, safety, and belonging.

\reviewsection{What the Unknowns Require From the Team}

The completed companion contains many ``unknown'' responses. That is an evidence boundary, not an excuse. One event cannot establish time patterns, sensory conditions, family circumstances, praise preferences, or a stable behavioral function. A fuller PTR process would require:

\begin{itemize}
  \item the learner's perspective through an accessible communication route;
  \item family and team knowledge about familiar tools and successful transitions;
  \item repeated observations across routines and adults;
  \item review of existing IEP, communication, sensory, and health supports;
  \item data on transition latency, distress, participation, return, and recovery;
  \item documentation of which supports were offered and whether they were accessible; and
  \item adult implementation-fidelity data.
\end{itemize}

If the learner does not benefit, the team should reconsider the hypothesis, access route, or implementation before concluding that the student failed. Data must examine adult behavior and environmental design alongside student behavior.

\reviewsection{How PTR Will Inform My Future Classroom}

PTR gives me a structure for turning ``behavior is communication'' from a slogan into an accountable process. In my future classroom, I will:

\begin{enumerate}
  \item define behavior in observable language;
  \item separate what I saw from what I inferred;
  \item collect information about tasks, transitions, tools, sensory conditions, peers, communication access, and adult responses;
  \item ask students and families what the behavior, support, and goal mean to them;
  \item prevent unnecessary barriers through visible and flexible lesson design;
  \item teach a functionally useful communication or participation route;
  \item reinforce agency and meaningful engagement;
  \item monitor student outcomes and adult implementation; and
  \item revise a plan that reduces visible behavior without increasing dignity, safety, communication, or belonging.
\end{enumerate}

\reviewsection{Puzzle Plan and EQUITAS Audit}

The Puzzle Plan is my interdisciplinary framework for combining lived experience, education, data, and research to understand complex problems more fully. EQUITAS is the equity-focused lens within that framework; it asks whose knowledge is included, what context is missing, and where power or bias may shape decisions.

Applied to PTR, this means no behavior, diagnosis, checklist response, or adult interpretation should stand alone. The learner's protest is one piece of evidence, and the adult's transition design is another. Keeping a familiar tool may be an access support rather than a reward for ``bad behavior.'' Communication modes should not be ranked by how typical they appear, and a learner should not have to surrender regulation or identity to be counted as participating.

\reviewsection{Conclusion}

Completing the PTR FBA Checklist changed the question from \textit{``How do I stop the protest?''} to \textit{``What happened around this transition, what might the learner be communicating, and what changes reveal participation?''} The most important evidence in this case is that taking away the laptop was not required for engagement. When the unnecessary either-or condition changed, the learner joined the blocks activity.

That result does not prove one function, and it does not create a universal rule about laptops. It demonstrates why support planning must remain collaborative, testable, and humble. A positive plan should prevent avoidable barriers, teach communication that works, reinforce meaningful access, and revise itself when the student's life does not improve. PTR is most useful when it makes educators more accountable for the environments and responses they design, not merely more precise in describing a student's behavior.

\small
\bibliography{references}

\end{document}
